\section{The observation and the main results}
\label{sec:setting}
The number of outcomes of a sensor and the information in its controls
are different resources. Here every preparation returns one bit, but
the controller chooses its spatial input. Our question is how many such
preparations are needed for geometric reconstruction, once their input
precision is specified. The reciprocal identity below reduces collision
observations to an occupation difference. The main point is that this
reduction can be evaluated at isolated, adaptively chosen locations;
a fine grid throughout the observation aperture is unnecessary.
We distinguish a localized experiment with worst-case implementation
errors from a stationary-noise experiment with a fixed calibrated launch
support. The latter, introduced below, has no shrinking random launch
spread and does not require the launch density to be known.

\subsection{Physical class and loss}
Fix $0<\beta\le1$ and put $s=6+\beta$. A table is the nonempty union
\begin{equation}\label{eq:presentation}
 \mathcal O=\bigcup_{i=1}^r(C_i+L\Z^2),\qquad 1\le r\le r_0.
\end{equation}
This presentation has exactly $r$ component orbits. It need not be
primitive. The constants in the following bounds are fixed and known:
\begin{equation}\label{eq:priors}
 \|L\|,\|L^{-1}\|\le L_0,\quad |\det L|\ge V_0>0,\quad
 \diam C_i\le D_0,\quad \dist(C,D)\ge d_0>0
\end{equation}
for distinct physical components. All components are strictly convex,
with curvature in $[\kappa_-,\kappa_+]\subset(0,\infty)$ and centered
support functions bounded by $M$ in $C^{6,\beta}(\R/2\pi\Z)$.
Reflections, repeated shapes and nonprimitive presentations are allowed.
Write $\mathcal B$ for this class, $\Pi=\{w:\mathcal O+w=\mathcal O\}$
for its full period group, $z_C$ for a component's Steiner point, and
$p_C$ for its centered support function. All norms on finite-dimensional
spaces are fixed once and for all.

A known positive nonperiod-patch margin $\eta$ will be required only
for uniform finite discrete decisions. Its precise definition is
\eqref{eq:patch-margin}. We denote the resulting subclass by
$\mathcal B_\eta$. The geometric loss is the maximum $C^2$ error of
matched component support functions in a protected laboratory patch,
together with the error of a matched primitive period basis and free
area. Representatives are chosen in a common protected patch. The
basis is matched to a true basis, not required to be a continuously
chosen reduced basis. Exact rational relations and the primitive orbit
count are stated separately from this real-valued loss.

\subsection{Commands and their accounting}
For a nominal displacement $a$, let $B_a(q)$ equal one if a free start
at $q$ has a first collision on the unreflected segment $[q,q+a]$.
It equals zero for a solid start and for a free miss. Both are attempted
preparations and enter the denominator. For pointwise statements solids
and swept sets are closed: a boundary solid start is zero, and a start
on a swept boundary but outside every solid is one. This fixes the
tangency convention; changing it on null sets has no effect on the means.
Fix an even smooth probability kernel $\kappa$ supported in the unit disk,
and set
\[
 f_{x,\sigma}(q)=\sigma^{-2}\kappa((q-x)/\sigma),\qquad
 u_\sigma=\kappa_\sigma*\one_\mathcal O.
\]
Choose $\sigma_0,t>0$ with $t+2\sigma_0<d_0$, and prescribe the compass
law
\begin{equation}\label{eq:compass}
 \rho=\tfrac14(\delta_{te_1}+\delta_{-te_1}
                         +\delta_{te_2}+\delta_{-te_2}).
\end{equation}
A forward command draws $q\sim f_{x,\sigma}$ and $a\sim\rho$
independently and attempts $(q,a)$. A reverse command draws that same
nominal joint law and attempts $(q+a,-a).$ Separate fresh preparations
are used; pairing the same physical particle is unnecessary. Only the
two pooled means are used, not means resolved by direction. The joint
reverse law, rather than only its marginal direction law, is prescribed.

In this localized bounded-error experiment, each actual preparation may be coupled to its nominal one with position,
duration and angle errors bounded by $\ell,\tau,\alpha$, respectively.
It remains a unit-speed straight flight until its first impact. The
bounds hold conditionally on past commands, and preparations at a given
command are fresh and conditionally independent. Section~\ref{sec:queries}
proves the bias bound $C(\ell+\tau+t\alpha)/\sigma$. Calibration is a
resource of the controller, not an estimate inferred from the bits.
The launch aperture is fixed by the prior and does not grow as accuracy
increases. It includes the translated reverse starts and the kernel
supports. No free point, obstacle label, impact location, impact angle,
collision time, primitive cell or area normalization is supplied.

\begin{theorem}[Adaptive scalar reconstruction]\label{thm:adaptive}
Fix the numerical bounds of $\mathcal B_\eta$ and the kernel above.
There are $C,c,\nu_0>0$ such that, for $0<\nu<\nu_0$ and
$0<\delta<1/4$, a finite adaptive experiment with the reciprocal
compass commands reconstructs a smooth strictly convex periodic
presentation with geometric loss at most $C\nu$, with probability at
least $1-\delta$. It recovers the true primitive orbit count and exact
rational relations of a true period-generating list in a true independent-pair basis. The Euclidean coordinates of these vectors are estimates.

After a prior-dependent coarse acquisition, the finest localization
scale and sufficient coupled calibration are
\begin{equation}\label{eq:precision}
 \sigma=c\nu^{s/(s-2)},\qquad \ell+\tau+t\alpha\le c\sigma.
\end{equation}
The number of nominal command centers, counted with repetitions if
necessary, and the number of attempted preparations satisfy
\begin{align}
 J_\nu&\le C\nu^{-1/(s-2)}\log(C/\nu),\label{eq:centers}\
 N_\nu&\le C\nu^{-1/(s-2)}\log(C/\nu)
                               \log(C/(\nu\delta)).\label{eq:attempts}
\end{align}
Every solid start and free miss is charged in $N_\nu$. All commands lie
in one fixed prior-controlled aperture. The reconstruction uses finite
local computations, radial bisections and finite polynomial data, not
a search over an infinite-dimensional class of physical candidates.
Its constants include the exit-time and patch-margin bounds.
\end{theorem}

Here and below a finite smooth partition of unity, with its derivatives,
is part of the fixed numerical representation. Certified evaluations
to the reserved tolerances suffice. The theorem counts preparations and
centers, not apparatus travel or bit operations for elementary functions.
Known $s$ is used; adaptivity refers to the locations of commands, not
adaptation to unknown smoothness.

\begin{theorem}[A necessary number of binary observations]
\label{thm:lower}
For each fixed $0<\beta\le1$ there is a fixed choice of the bounds
above and a nondegenerate subclass $\mathcal P\subset\mathcal B_\eta$
with the following property. Any adaptive randomized procedure whose
only table-dependent outputs are at most $N$ bits, and which recovers
the component boundaries with $C^2$ error at most $\nu$ with probability
at least $3/4$ uniformly on $\mathcal P$, must satisfy
\begin{equation}\label{eq:lower}
                         N\ge c\nu^{-1/(s-2)}.
\end{equation}
This holds even when the lattice, its primitive orbit count, component
correspondences and laboratory pose are given, and implementation is
exact. Input choices may have arbitrary precision and depend on all
previous bits. Independent controller randomness is allowed.
\end{theorem}

The construction of $\mathcal P$ is given in Section~\ref{sec:lower}.
The assertion is not a lower bound for every subclass, in particular
not for a singleton or a finite parametric family. On any bounded class containing this construction, $\eqref{eq:attempts}$ and
$\eqref{eq:lower}$ identify the power of the attempted-bit complexity.
A logarithmic gap remains. Neither result prices input precision or
converts an active experiment into a passive billiard invariant.
